\section*{第 \ref{ch:b1:genomes} 章——细胞与病毒的基因组}

\begin{solution}{exo:b1:genomes:1}
细菌：一条环状\omterm{def:b1:genomes:genome}{染色体}（外加\omterm{def:b1:genomes:genome}{质粒}），位于
无膜的\omterm{def:b1:genomes:genome}{拟核}中，借助超螺旋和少量碱性
蛋白压缩，约 \qty{90}{\%} 为编码序列。真核生物：\omterm{def:b1:eukaryotic-cell:nucleus}{细胞核}中有多条线性
\omterm{def:b1:genomes:genome}{染色体}，缠绕\omterm{def:b1:genomes:nucleosome}{组蛋白}形成\omterm{def:b1:genomes:nucleosome}{染色质}，编码序列
仅占百分之几。
\end{solution}

\begin{solution}{exo:b1:genomes:2}
147 个\omterm{thm:b1:nucleic-acids:helix}{碱基对}缠绕在\omterm{def:b1:genomes:nucleosome}{组蛋白}八聚体上约 1.7 圈，八聚体
由 H2A、H2B、H3、H4 各两个组成，由 H1 封住，每 200 对碱基一个。串珠结构
$\times 6$；\qty{30}{nm} 纤维 $\times 40$；环 $\times 1000$；
中期\omterm{def:b1:genomes:genome}{染色体} $\times \num{10000}$。
\end{solution}

\begin{solution}{exo:b1:genomes:3}
\qty{1.5}{\%}；转座元件及其残迹，占 \qty{45}{\%}。
\end{solution}

\begin{solution}{exo:b1:genomes:4}
\omterm{def:b1:genomes:genome}{基因组}（\omterm{def:b1:nucleic-acids:chain}{DNA} 或 \omterm{def:b1:nucleic-acids:chain}{RNA}）装在\omterm{def:b1:proteins:peptide}{蛋白质}\omterm{def:b1:genomes:virus}{衣壳}内，有时还有包膜，
只能在\omterm{def:b1:cell-unit-of-life:cell}{细胞}内借助\omterm{def:b1:cell-unit-of-life:cell}{细胞}的机器复制。离开\omterm{def:b1:cell-unit-of-life:cell}{细胞}后，它没有代谢、不生长、
无应答——是一个惰性颗粒；它之所以“活”，
只在于携带可遗传并会演化的
信息。
\end{solution}

\begin{solution}{exo:b1:genomes:5}
$\num{4.6e6}\times 0.34 = \qty{1.56}{mm}$；$\num{4.6e6}/200 = \num{23000}$
个\omterm{def:b1:genomes:nucleosome}{核小体}。细菌实际采用的方式是：旋转酶引入负超螺旋，\omterm{def:b1:nucleic-acids:chain}{DNA} 折叠成
约五十个环并附着于\omterm{def:b1:proteins:peptide}{蛋白质}核心，少量碱性蛋白使 \omterm{def:b1:nucleic-acids:chain}{DNA} 弯曲，
压缩约千倍——形成一微米大小的\omterm{def:b1:genomes:genome}{拟核}。
\end{solution}

\begin{solution}{exo:b1:genomes:6}
\omterm{def:b1:genomes:gene}{外显子} $1500/8 = \qty{190}{bp}$；\omterm{def:b1:genomes:gene}{内含子} $25\,500/7 = \qty{3640}{bp}$；
$25.5/27 = \qty{94}{\%}$ 的转录本被剔除。
\end{solution}

\begin{solution}{exo:b1:genomes:7}
\emph{E. coli} 每 Mb 930 个\omterm{def:b1:genomes:gene}{基因}；酵母 500；果蝇 100；人类 6；肺鱼
0.15。\omterm{def:b1:genomes:gene}{基因}密度下降了四个数量级，而\omterm{def:b1:genomes:gene}{基因}数量
增加了五倍：多出的 \omterm{def:b1:nucleic-acids:chain}{DNA} 是非\omterm{def:b1:genomes:gene}{基因}的。
\end{solution}

\begin{solution}{exo:b1:genomes:8}
$10^6\times 300 = \qty{3e8}{bp}$：占\omterm{def:b1:genomes:genome}{基因组}的 \qty{9}{\%}。$60\times
10^6/20 = \num{3e6}$ 代，平均每三代在谱系中的某处出现
一次新的插入。
\end{solution}

\begin{solution}{exo:b1:genomes:9}
卵提供\omterm{def:b1:cell-unit-of-life:cell}{细胞质}及其中的\omterm{def:b1:eukaryotic-cell:mitochondrion}{线粒体}；精子的少量\omterm{def:b1:eukaryotic-cell:mitochondrion}{线粒体}
在受精后被排除或降解。因此，\omterm{def:b1:eukaryotic-cell:mitochondrion}{线粒体} \omterm{def:b1:nucleic-acids:chain}{DNA} 不经混合地
由母亲传给子代，其累积的突变可以
追溯母系谱系的历史。
\end{solution}

\begin{solution}{exo:b1:genomes:10}
营养充足的培养物，两小时：溶原途径得到 $2^6 = 64$ 个拷贝（每次宿主分裂得一个）；
裂解途径在 \qty{40}{min} 内产生 100 个噬菌体，
若宿主充足，其后代两小时内可达 $100^3 = 10^6$ 个——裂解途径远占优势。
饥饿的培养物：溶原途径只有一个拷贝，安全地藏在活\omterm{def:b1:cell-unit-of-life:cell}{细胞}内；裂解途径产生 100 个噬菌体，
却没有可以侵染的\omterm{def:b1:cell-unit-of-life:cell}{细胞}，它们会失活。开关
读取宿主状态，是因为最优策略取决于
今后是否还会有新宿主可用。
\end{solution}

\begin{solution}{exo:b1:genomes:11}
$\num{1.3e11}\times 0.34\,\mathrm{nm} = \qty{44}{m}$（单倍体）；
$\num{6.5e8}$ 个\omterm{def:b1:genomes:nucleosome}{核小体}；\omterm{def:b1:genomes:nucleosome}{组蛋白} $\num{6.5e8}\times 108\,000\times
1.66\times 10^{-24} = \qty{117}{pg}$。复制起点：$\num{1.3e6}$ 个，每个
复制起点向两侧各复制 \qty{50}{kb}，速率 \qty{50}{bp/s}：\qty{1000}{s}
——时间不是问题，只要复制起点足够多；代价是
\omterm{def:b1:nucleic-acids:nucleotide}{核苷酸}（每次分裂需 \qty{130}{pg} 的 \omterm{def:b1:nucleic-acids:chain}{DNA}）、\omterm{def:b1:genomes:nucleosome}{组蛋白}、
合成它们所需的时间，以及大得多的\omterm{def:b1:eukaryotic-cell:nucleus}{细胞核}和\omterm{def:b1:cell-unit-of-life:cell}{细胞}，
这正是此类\omterm{def:b1:cell-unit-of-life:cell}{细胞}分裂缓慢的原因。
\end{solution}

\begin{solution}{exo:b1:genomes:12}
设计只包含所需之物，不多一分；记录则
保存已经发生的一切。假\omterm{def:b1:genomes:gene}{基因}是曾被复制、后因突变而失效的\omterm{def:b1:genomes:gene}{基因}的遗骸；
转座元件是寄生者，它们自我复制了数千万年，
从未被清除；\omterm{def:b1:genomes:gene}{基因}家族显示先复制、后分化；
C 值悖论表明，\omterm{def:b1:nucleic-acids:chain}{DNA} 的数量反映的是谱系积累与丢失的历史，
而非其需要。选择使\omterm{def:b1:genomes:gene}{基因}保持功能，
其余部分则任其漂变：我们在\omterm{def:b1:genomes:genome}{基因组}中
读到的，大多是
历史。
\end{solution}

\begin{solution}{pb:b1:genomes:1}
\textbf{1.} $\num{6.4e9}\times 0.34\,\mathrm{nm} = \qty{2.18}{m}$。
\textbf{2.} $\num{6.4e9}/10 = \num{6.4e8}$ 圈。
\textbf{3.} $\num{2.5e8}\times 0.34 = \qty{8.5}{cm}$；\qty{1.7}{cm}。
\textbf{4.} $\frac{4}{3}\pi\times 3^3 = \qty{113}{\micro m^3}$。
\textbf{5.} $\pi\times(1\,\mathrm{nm})^2\times 2.18\,\mathrm{m} =
\qty{6.8e-18}{m^3} = \qty{6.8}{\micro m^3}$：占\omterm{def:b1:eukaryotic-cell:nucleus}{细胞核}的 \qty{6}{\%}。
\textbf{6.} $\num{6.4e9}\times 650\times 1.66\times 10^{-24} =
\qty{6.9e-12}{g} = \qty{6.9}{pg}$。
\textbf{7.} \omterm{def:b1:eukaryotic-cell:nucleus}{细胞核}质量 $113\times 10^{-12}\,\mathrm{mL}\times 1.1 =
\qty{124}{pg}$；\omterm{def:b1:nucleic-acids:chain}{DNA} 占其 \qty{5.6}{\%}，\omterm{def:b1:proteins:peptide}{蛋白质}占 \qty{20}{\%}。
\textbf{8.} $\num{6.4e9}/200 = \num{3.2e7}$。
\textbf{9.} $\num{3.2e7}\times 108\,000\times 1.66\times 10^{-24} =
\qty{5.7}{pg}$：与 \omterm{def:b1:nucleic-acids:chain}{DNA} 量大致相当。
\textbf{10.} $68/11 = 6.2$。
\textbf{11.} 六个\omterm{def:b1:genomes:nucleosome}{核小体}，\qty{1200}{bp} $= \qty{408}{nm}$ 的螺旋，
压入 \qty{11}{nm} 之内：倍数 $37$。
\textbf{12.} $2.18/37 = \qty{5.9}{cm}$ 的纤维——相当于
\omterm{def:b1:eukaryotic-cell:nucleus}{细胞核}直径的一万倍，折叠在\omterm{def:b1:eukaryotic-cell:nucleus}{细胞核}之内。
\textbf{13.} 全部 $\num{3.2e7}$ 个，另加第二条拷贝上
同样数目的新\omterm{def:b1:genomes:nucleosome}{核小体}：共 $\num{6.4e7}$ 个，历时 \qty{28800}{s}，
约每秒 \num{2200} 个。
\textbf{14.} 每个环含纤维 $75\,000\times 0.34/37 = \qty{690}{nm}$；
共 $\num{6.4e9}/75\,000 = \num{85000}$ 个环。
\textbf{15.} $\qty{8.5}{cm}/\qty{10}{\micro m} = 8500$。
\textbf{16.} $\pi\times 0.35^2\times 10 = \qty{3.85}{\micro m^3}$；\omterm{def:b1:nucleic-acids:chain}{DNA}
$\pi\times 10^{-6}\times 0.085\,\mathrm{m} = \qty{0.27}{\micro m^3}$：
\qty{7}{\%}——其余为\omterm{def:b1:genomes:nucleosome}{组蛋白}、支架蛋白和水。
\textbf{17.} \omterm{def:b1:nucleic-acids:chain}{DNA} 总量为 $\num{6.4e9}$ 对，密度相同
（每条染色单体中 $\num{2.5e8}$ 对位于 \qty{3.85}{\micro m^3} 内；每条\omterm{def:b1:genomes:genome}{染色体}
两条染色单体）：$2\times 6.4\times 10^9/2.5\times 10^8\times
3.85 = \qty{197}{\micro m^3}$，比间期\omterm{def:b1:eukaryotic-cell:nucleus}{细胞核}更大，
因为此时的\omterm{def:b1:genomes:genome}{染色体}是彼此间隔的棒状结构。
\textbf{18.} 间期\omterm{def:b1:genomes:nucleosome}{染色质}必须能被读取和复制：聚合酶
及其辅助因子需要接近 \omterm{def:b1:nucleic-acids:chain}{DNA}，因此大部分保持为
开放的 \qty{30}{nm} 纤维和环；中期\omterm{def:b1:genomes:nucleosome}{染色质}是惰性的，可以
被压缩以便运输。
\textbf{19.} $\num{6.4e9}$ 个位点，每秒 $10^3$ 个：$\num{6.4e6}$
秒，合 74 天。\omterm{def:b1:proteins:peptide}{蛋白质}并非随机搜索：它们非特异地结合 \omterm{def:b1:nucleic-acids:chain}{DNA}
并沿其滑动，许多拷贝并行搜索，
因此几秒内即可找到一个位点。
\textbf{20.} $\num{20000}\times 27\,\mathrm{kb} = \qty{540}{Mb}$，
占 \qty{3.2}{Gb} 的 \qty{17}{\%}；编码区 $\num{20000}\times 1.5 =
\qty{30}{Mb}$，\qty{0.9}{\%}（若较短\omterm{def:b1:genomes:gene}{基因}计入得更仔细，约为 1.5\%
）。
\textbf{21.} $113\times 2\times 250/6400 = \qty{8.8}{\micro m^3}$（两份
拷贝）；$\num{5e8}/8.8 = \qty{5.7e7}{bp/\micro m^3}$，与
整个\omterm{def:b1:eukaryotic-cell:nucleus}{细胞核}相同。
\textbf{22.} $\num{2.6e11}/\num{5.7e7} = \qty{4600}{\micro m^3}$：
直径 \qty{21}{\micro m}，体积为人类\omterm{def:b1:eukaryotic-cell:nucleus}{细胞核}的四十倍。
\textbf{23.} $\num{3.2e7}$ 个八聚体，$\num{3.2e10}$ 个残基——相当于
\omterm{def:b1:cell-unit-of-life:cell}{细胞}每日\omterm{def:b1:proteins:peptide}{蛋白质}合成量的十分之一，全部用于包装。
\textbf{24.} 细菌的\omterm{def:b1:genomes:genome}{染色体}短约千倍，环状超螺旋
足以将其装入一微米之内；真核生物
必须把一段长度为其\omterm{def:b1:eukaryotic-cell:nucleus}{细胞核}一万倍的 \omterm{def:b1:nucleic-acids:chain}{DNA} 折叠起来，
并需要一套多级包装体系，其中\omterm{def:b1:genomes:nucleosome}{核小体}是第一级。
\textbf{25.} 约 \num{8500} 倍（\qty{8.5}{cm} 压入 \qty{10}{\micro m}），
其构成为：$\times 6$（\omterm{def:b1:genomes:nucleosome}{核小体}），再 $\times 6$（\qty{30}{nm} 纤维，
共 $\times 37$），再 $\times 27$（环，共 $\times 1000$），
以及再 $\times 8$（中期凝缩）。
\end{solution}
